% Unicode in the archived scientific text, represented using standard math glyphs.
\newunicodechar{α}{\ensuremath{\alpha}}
\newunicodechar{λ}{\ensuremath{\lambda}}
\newunicodechar{η}{\ensuremath{\eta}}
\newunicodechar{θ}{\ensuremath{\theta}}
\newunicodechar{ℓ}{\ensuremath{\ell}}
\newunicodechar{×}{\ensuremath{\times}}
\newunicodechar{≥}{\ensuremath{\geq}}
\newunicodechar{≤}{\ensuremath{\leq}}
\newunicodechar{±}{\ensuremath{\pm}}
\newunicodechar{≈}{\ensuremath{\approx}}
\newunicodechar{→}{\ensuremath{\rightarrow}}
\newunicodechar{√}{\ensuremath{\surd}}
\newunicodechar{‖}{\ensuremath{\Vert}}
\newunicodechar{₀}{\ensuremath{{}_0}}
\newunicodechar{⁻}{\ensuremath{{}^{-}}}
\newunicodechar{¹}{\ensuremath{{}^{1}}}
\newunicodechar{²}{\ensuremath{{}^{2}}}
\newunicodechar{³}{\ensuremath{{}^{3}}}
\newunicodechar{⁴}{\ensuremath{{}^{4}}}
\newunicodechar{⁵}{\ensuremath{{}^{5}}}
\newunicodechar{⁶}{\ensuremath{{}^{6}}}
\newunicodechar{⁷}{\ensuremath{{}^{7}}}
\newunicodechar{ᵀ}{\ensuremath{{}^{\mathsf T}}}
